\begin{proof}[Proof of Proposition~\ref{sga3:XXI:3.3.6}]
Choose \(w\in W\) such that
\(\operatorname{Card}(w(R_+)\cap P')\) is maximal. Replacing \(R_+\)
by \(w(R_+)\), we may assume that
\begin{equation}
 \operatorname{Card}(R_+\cap P')
 \geq
 \operatorname{Card}\bigl(R_+\cap s_\alpha(P')\bigr)
 \quad\text{for every }\alpha\in\mathscr{S}(R_+)=\Delta.
 \tag{\(\ast\)}\label{sga3:XXI:3.3.6:star}
\end{equation}
We show that \(R_+\subset P'\). Otherwise, since \(P'\) is closed,
there is an \(\alpha\in\Delta\) such that \(\alpha\notin P'\). As
\(P'\) satisfies (P 2), one then has \(-\alpha\in P'\), and also
\(-2\alpha\in P'\) if \(2\alpha\) is a root.

For every \(\beta\in R_+\cap P'\), one has \(\beta\neq\alpha\). If
\(2\alpha\) is not a root, then
\(s_\alpha(\beta)\in R_+\) by Lemma~\ref{sga3:XXI:3.3.1}; hence
\[
 s_\alpha(R_+\cap P')
 \subset R_+\cap s_\alpha(P').
\]
But \(R_+\cap s_\alpha(P')\) also contains
\(\alpha=s_\alpha(-\alpha)\), contradicting
\eqref{sga3:XXI:3.3.6:star}. If \(2\alpha\) is a root, the same
argument applies.
\end{proof}

Thus, to study closed sets of roots satisfying (P 2), it is enough to
consider those that contain a system of positive roots.

\begin{proposition}\label{sga3:XXI:3.3.10}
Let \(R_+\) be a system of positive roots, and let \(P'\) be a closed
subset of \(R\) containing \(R_+\). Put
\[
 \Delta=\mathscr{S}(R_+),
 \qquad
 \Delta'=\Delta\cap(-P').
\]
Then \(P'\) is the union of \(R_+\) and the set of roots that are
linear combinations of the elements of \(\Delta'\) with negative
coefficients.
\end{proposition}

\begin{proof}
We argue by induction on the order of a root
\(\gamma\in P'\setminus R_+\). If
\(\operatorname{ord}_\Delta(\gamma)=-1\), then
\(-\gamma\in\Delta\), and hence \(\gamma\in-\Delta'\).

Suppose that \(\operatorname{ord}_\Delta(\gamma)<-1\). By
Proposition~\ref{sga3:XXI:3.1.2}, there exists \(\beta\in\Delta\)
such that \(-\gamma-\beta\in R\). Then
\[
 0<\operatorname{ord}_\Delta(-\gamma-\beta)
  =\operatorname{ord}_\Delta(-\gamma)-1
  <\operatorname{ord}_\Delta(-\gamma).
\]
The first inequality shows that \(-\gamma-\beta\in R_+\).
Consequently, since \(R_+\cap(-R_+)=\varnothing\), and since
\(\gamma+\beta\) is the sum of two roots of \(P'\), the root
\(\gamma+\beta\) belongs to \(P'\setminus R_+\) and satisfies
\[
 \operatorname{ord}_\Delta(\gamma+\beta)
 >\operatorname{ord}_\Delta(\gamma).
\]
By induction, \(\gamma+\beta\) is a linear combination of the
elements of \(\Delta'\) with negative coefficients. Since
\(\gamma=(\gamma+\beta)-\beta\), it remains only to show that
\(\beta\in-P'\). But \(\gamma\in P'\) and
\(-(\gamma+\beta)\in R_+\subset P'\), so
\[
 -\beta=\gamma-(\gamma+\beta)\in P'.
\]
\end{proof}

\begingroup
\renewcommand{\thetheorem}{3.3.10.1}
\begin{definition}\label{sga3:XXI:3.3.10.1}
\footnote{Editors' note: This definition was inserted here.}
A subset \(R'\) of \(R\) is called \emph{symmetric} if
\[
 -R'=R'.
\]
\end{definition}
\endgroup
\addtocounter{theorem}{-1}
